\begin{abstract}
\noindent
Large language models are increasingly proposed to assist peer review, yet most evaluations
judge the prose of machine-generated review \emph{text}, not the validity of the numeric
\emph{score} a system assigns. We validate \textsc{AIPR}, which emits a weighted overall
manuscript score, against the public decision outcomes of a major machine learning venue,
grading by prompting alone. With hypotheses pre-registered before any score met any outcome, the
production score separates rejected from accepted submissions (AUROC~\aurocFull), replicated on a
\NminiPrimary-paper cohort (\aurocMini); it rises across tiers, tracks the mean reviewer rating,
and is strongest where we claim it, with the lowest-scoring fifth rejected far above the base
rate. The validity comes mostly from the model: a one-paragraph prompt discriminates almost as
well (the small gap does not meet the pre-declared criterion, \aurocDiffPrel). What the
engineering adds is reliability: AIPR's score barely moves across runs (within-paper SD
\fullRunSD{} vs.\ \naiveRunSD{} points, \relPairedP{}), a gap a balanced four-arm follow-up
confirms at both model tiers and finds largest on rejects, together with a grounded review.
\end{abstract}
